\section{Safeguarding LLM Agents from Misalignment through Provenance Analysis}

Yining She, Yiliang Liang, Eunsuk Kang. \emph{Safeguarding LLM Agents from Misalignment
through Provenance Analysis.} ASE~2026, arXiv:2607.01236v2. Репликационный пакет:
\url{https://github.com/cmu-soda/ProvenanceGuard}, Zenodo 10.5281/zenodo.19342590.

\cite{she2026provenanceguard}

\subsection{Ключевые идеи, математическое описание метода}
Несогласованность (misalignment)~-- действие агента, отклоняющееся от намерения
пользователя, без злого умысла модели. Три типа: неверный инструмент, неверные параметры,
неверная интерпретация недоопределённого запроса. Ключевая идея: действие согласовано,
если оно \emph{обосновано прослеживаемыми уликами в контексте}; проверка сводится
к установлению отношений провенанса между элементами действия и элементами контекста.

Контекст guardrail'а: запрос $q$, документация инструментов $D$ и история
\begin{equation}
h = \langle ((t_1,a_1),o_1), \dots, ((t_k,a_k),o_k) \rangle,
\end{equation}
где $(t_i,a_i)$~-- вызов инструмента с назначением параметров, $o_i$~-- наблюдаемый вывод.
Guardrail~-- функция
\begin{equation}
\mathrm{guard}(t,a,q,D,h) \in \{\text{Aligned}, \text{Misaligned}\}.
\end{equation}
Пять отношений (табл.~1 статьи): $\mathrm{SpecifiedBy}(t,d)$, $\mathrm{ContributeTo}(s,q)$,
$\mathrm{RelevantTo}(d,s)$, $\mathrm{DerivedFrom}(p,v,e)$, $\mathrm{CanAddress}(d,a,s)$,
где $s$~-- подзадача из множества $S$, $d\in D$~-- документ инструмента, $e$~-- улика.
Схематическая запись трёх предикатов:
\begin{equation}
\begin{aligned}
&\mathrm{ToolProven}(t,q,D)\triangleq \exists d\in D,\ s\in S:\\
&\quad\mathrm{SpecifiedBy}(t,d)\wedge\mathrm{ContributeTo}(s,q)\wedge\mathrm{RelevantTo}(d,s).
\end{aligned}
\end{equation}
\begin{equation}
\begin{aligned}
&\mathrm{ParamProven}(a,d_t,s_t,q,h)\triangleq \mathrm{CanAddress}(d_t,a,s_t)\ \wedge\\
&\quad\forall (p,v)\in a\ \exists e\in E(q,h):\ \mathrm{DerivedFrom}(p,v,e).
\end{aligned}
\end{equation}
\begin{equation}
\begin{aligned}
&\mathrm{InterpProven}(t,a,q,D,h,s_t)\triangleq \neg\exists (t',a'):\\
&\quad(t',a')\not\equiv(t,a)\ \wedge\ \mathrm{Admissible}(t',a',q,D,h,s_t).
\end{aligned}
\end{equation}
Здесь Admissible сокращённо обозначает допустимость альтернативы для той же подзадачи
по проверкам инструмента и параметров; $\equiv$~-- эквивалентность действий.
Множество улик $E(q,h)$~-- факты, значения и ограничения, явно присутствующие в $q$
или в выводах $o_i$ предыдущих вызовов. DerivedFrom не требует точного совпадения:
значение $1700$ считается выведенным из улики <<два месяца по $850$>>. Параметры делятся
на \emph{выводимые} (должны иметь источник в контексте) и \emph{генеративные} (свободный
текст); проверяются только выводимые. В формуле (4) оригинала квантор записан по всему $a$,
но описание и реализация ограничивают проверку множеством $a_{der}\subseteq a$;
при $a_{der}=\varnothing$ этап II пропускается.


\subsection{Алгоритм}
Трёхэтапный конвейер с ранней остановкой; каждый этап отвечает за один тип несогласованности.
\begin{enumerate}
  \item Фильтр: инструмент должен существовать в $D$ (механически) и быть помечен
        разработчиком как меняющий среду (по метаданным); только читающие инструменты
        не проверяются.
  \item Этап I (инструмент): подзадача $s$ берётся из плана самого агента; LLM проверяет
        $\mathrm{RelevantTo}(d_t,s)$ и $\mathrm{ContributeTo}(s,q)$.
  \item Этап II (параметры): LLM определяет множество выводимых параметров $a_{der}\subseteq a$
        по документации $d_t$, затем для каждого $(p,v)\in a_{der}$ ищет улику $e$ и проверяет
        $\mathrm{DerivedFrom}$ и $\mathrm{CanAddress}$.
  \item Этап III (интерпретация): LLM, \emph{не видя предложенного действия}, проверяет,
        существует ли иное допустимое действие для $s_t$; если да, запрос недоопределён,
        действие блокируется.
  \item При провале любого этапа~-- вердикт Misaligned и запрос подтверждения у пользователя.
\end{enumerate}
Сравниваемые методы: DirectPrompt (LLM-as-a-judge одним промптом) и ProvePrompt
(те же три проверки одним промптом, без разделения на этапы).

\subsection{Наборы данных}
\begin{itemize}
  \item \emph{Agent-SafetyBench} (подмножество failure mode~2): 188 трасс сгенерированы
        GPT-5, вручную размечены двумя авторами ($\kappa=0{,}85$), 6 отброшены;
        итого 148 несогласованных и 34 согласованных трассы, в среднем 2,2 шага.
  \item \emph{WorkBench}: трассы GPT-4 из бенчмарка; несогласованная~-- хотя бы один вызов
        расходится с эталонной последовательностью действий (62), task-successful~-- 326.
        Успешная трасса не обязательно согласована (недоопределённость не размечена).
  \item План агента для обоих наборов восстановлен GPT-5.
\end{itemize}

\subsection{Метрики}
Для пары (метод, LLM) и $R=5$ повторов на трассу:
\begin{equation}
\mathrm{ER} = \frac{\sum_{r=1}^{R}\#\mathrm{FN}_r}{N\cdot R}, \qquad
\mathrm{IR} = \frac{\sum_{r=1}^{R}\#\mathrm{Intervene}_r}{N\cdot R},
\end{equation}
где ER считается на несогласованных трассах (пропуски), IR~-- на согласованных
(ложные вмешательства) или на task-successful (операционная нагрузка). Пять бинарных
исходов усредняются в оценку трассы $s^{(c)}\in\{0, 0{,}2,\dots,1\}$; статистика~-- средняя
парная разность $\Delta=\frac1N\sum_c (s^{(c)}_{\mathrm{PG}}-s^{(c)}_{m})$, односторонний
перестановочный тест со сменой знака (10\,000 перестановок, Монте-Карло), поправка
Холма по 11 LLM внутри каждой пары методов. Доверительные интервалы не приводятся.
Результаты: ER на Agent-SafetyBench $44{,}3\%\to14{,}7\%\to2{,}1\%$
(DirectPrompt~$\to$ ProvePrompt~$\to$ ProvenanceGuard), на WorkBench $32{,}4\%\to20{,}9\%\to18{,}7\%$;
IR на task-successful WorkBench $31{,}2\%\to13{,}0\%$. Стоимость: $\approx3{,}6$ вызова LLM
и 3\,979 входных токенов на случай против 1\,303 у DirectPrompt.
Ограничения: оценка проводится на уровне трасс; отдельные связи не получают
независимого эталона. Малая выборка согласованных случаев ограничивает выводы об IR;
отсутствие статистической значимости не доказывает отсутствие ухудшения.

\subsection{Модели}
11 backbone-моделей: GPT-5, GPT-5-mini, GPT-4.1, GPT-4.1-mini, Gemini-3-flash-preview,
Gemini-3.1-flash-lite-preview, Gemini-2.5-flash, Gemini-2.5-flash-lite, Claude-sonnet-4.6,
Claude-haiku-4.5, Qwen-3.5-9B (локально, vLLM). Параметры генерации по умолчанию.
